% Cleaned OCR chunk: chunks/09_shen_brush_talks.md
% Scope: Shen Kuo, Mengxi Bitan excerpts only.

\providecommand{\preservedimage}[1]{%
  \begin{figure}[htbp]
    \centering
    \IfFileExists{#1}{%
      \includegraphics[width=0.92\linewidth]{#1}%
    }{%
      \fbox{\parbox{0.84\linewidth}{\centering 圖像路徑保留：\texttt{#1}}}%
    }%
  \end{figure}%
}

\section*{夢溪筆談選段}

\subsection*{卷十八　技藝}

\subsubsection*{圍棋局數}

\paragraph{英譯補譯}
說書人相傳，唐代僧人一行曾經推算圍棋所有可能局面的總數，並認為可以窮盡。我細想之後，覺得此事本不難；只是數量極大，已非世間常用的數名可以表示。姑舉其大數：二路方盤、四枚棋子，可成八十一種局面；三路方盤、九枚棋子，可成一萬九千六百八十三種；四路方盤、十六枚棋子，可成四千三百零四萬六千七百二十一種；五路方盤、二十五枚棋子，可成八千四百七十二億八千八百六十萬九千四百四十三種。七路以上，數量已大到沒有通行名稱可記。至於十九路全盤三百六十一點，局面總數約為「萬」字連書四十三次的大數。

\subsubsection*{活版印刷}

\paragraph{英譯補譯}
慶曆年間，平民畢昇創製活版印刷。他的方法是用膠泥刻字，字模薄如銅錢邊緣，每一字各成一枚泥活字，再用火燒使之堅硬。先備一塊鐵板，在板上鋪松脂、蠟和紙灰之類的混合藥料。要印刷時，把鐵範放在鐵板上，在範中密排活字；排滿一範，就成為一塊穩固的印版。隨後近火烘烤，等背面的藥料稍稍熔化，再用平板壓平字面，使整版如磨石般平整。

若只印兩三本，這方法並不算簡便；若印數百、數千本，就極為神速。通常備兩版交替使用：一版正在刷印，另一版便同時排字；前一版印完，後一版也已備好，如此輪換，印刷便十分迅捷。

每個字都有數枚活字；常用字則各備二十枚以上，以應付同一頁中重複出現的字。活字不用時，按韻部分類，用紙簽標明，收在木格中。若遇到平日未備的生僻字，就臨時刻製，用草火燒成，片刻即可使用。

不用木活字，是因為木紋有粗有細，遇水後又會高低不平；排版時還會黏在藥料上，不易取下。因此不如用燒成的陶土活字。印完後，把印版再靠近火邊，使藥料熔化，用手一拂，活字便自行脫落，絲毫不被弄髒。

畢昇去世後，他的整副活字由沈括的族侄輩收藏，至當時仍視為珍寶。

\subsection*{卷二十一　異事（異疾附）}

\subsubsection*{虹}

\paragraph{原文}
世傳虹能入溪澗飲水，信然。

熙寧中，予使契丹，至其極北黑水境永安山下卓帳。是時新雨霽，見虹下帳前澗中。予與同職扣澗觀之，虹兩頭皆垂澗中。使人過澗，隔虹對立，相去數丈，中間如隔綃縠。自西望東則見；蓋夕虹也。立澗之東西望，則為日所鑠，都無所覩。久之稍稍正東，踰山而去。次日行一程，又復見之。孫彥先云：「虹乃雨中日影也，日照雨則有之。」

\paragraph{今譯}
世人傳說彩虹能進入溪澗飲水，這話確有其事。

熙寧年間，我出使契丹，到達它極北方的黑水境，在永安山下扎營。當時雨後剛剛放晴，看見一道彩虹垂到帳前的澗水中。我和同僚走近澗邊觀看，彩虹兩端都垂入澗水。又讓人渡過溪澗，隔著彩虹與我們相對而立，相距數丈，中間好像隔著薄紗。從西向東望去就能看見，因為那是傍晚的虹；若站在溪澗東邊向西看，便被日光照得耀眼，什麼也看不見。過了很久，彩虹漸漸向正東移動，越過山嶺而去。第二天走了一程，又再次看見它。孫彥先說：「虹是雨中的日影；日光照在雨點上，就會出現虹。」

\subsection*{卷二十四　雜誌一}

\subsubsection*{海陸變遷}

\paragraph{原文}
予奉使河北，遵太行而北，山崖之間，往往銜螺蚌殼及石子如鳥卵者，橫亙石壁如帶。此乃昔之海濱，今東距海已近千里。所謂大陸者，皆濁泥所湮耳。堯殛鯀於羽山，舊說在東海中，今乃在平陸。凡大河、漳水、滹沱、涿水、桑乾之類，悉是濁流。今關、陝以西，水行地中，不減百餘尺，其泥歲東流，皆為大陸之土，此理必然。

\paragraph{今譯}
我奉命出使河北，沿太行山向北而行，看見山崖之間常常夾有螺、蚌的殼，以及像鳥卵一樣的石子，橫貫石壁，形如帶狀地層。這裡從前應是海濱，如今向東距海已接近千里。所謂大陸，其實都是渾濁水流所挾帶的泥沙沉積、淤塞而成。堯誅殺鯀於羽山，舊說羽山在東海中，如今卻已在平陸。凡黃河、漳水、滹沱河、涿水、桑乾河一類，都是含沙量很高的濁流。如今潼關、陝州以西，河水下切地中，深處不下百餘尺；其中泥沙年年向東搬運，最終成為大陸之土，這在地貌變遷上是必然的道理。

\subsubsection*{指南針}

\paragraph{原文}
方家以磁石磨針鋒，則能指南，然常微偏東，不全南也。水浮多蕩搖。指爪及盌唇上皆可為之，運轉尤速，但堅滑易墜，不若縷懸為最善。其法取新纊中獨繭縷，以芥子許蠟，綴於針腰，無風處懸之，則針常指南。其中有磨而指北者。予家指南、北者皆有之。磁石之指南，猶柏之指西，莫可原其理。

\paragraph{今譯}
方術家用磁石摩擦針尖，針便能指南；但它常常稍微偏東，並不完全指向正南，這就是磁針的偏角現象。把磁針浮在水面上，常會飄蕩搖晃；放在指甲或碗口邊緣上也可以使它轉動，而且轉得更快，但這些支點堅硬光滑，磁針容易墜落，不如用絲縷懸掛最為穩妥。方法是從新絲綿中取一根獨繭絲，用芥子大小的一點蠟黏在針腰，把針懸在無風處，針便常常指南。其中也有摩擦後指北的。我家中指南、指北的磁針都有。磁石能指南，就像柏木枝幹趨西一樣，無法推究其原理。

\paragraph{原圖保留}
\preservedimage{images/d4a8886ce638fa3464f558b8c9a1932a570a557f66d4de747a34c69ade79d9e6.jpg}
\preservedimage{images/2f6c0487a50f0282edbf4c85b5d5c7d8454054cd248a8c5db183c048c711c25b.jpg}
\preservedimage{images/9ba1872eba396ba3fcb14170391f8c57060d494d202e9ea51e79d8c0ad1c9f3e.jpg}
\preservedimage{images/7b6c8c59583d67e076070ca60124f5b97f7e1f46d61eae8062b44497cf0ee5b2.jpg}
\preservedimage{images/a93d6ec79ea3faf441e39e04d4ffc9250f86fb9a33dc11d9399658e062718e4b.jpg}
\preservedimage{images/fbb24748342b96f2ec55f761c5b08597fd7647ff4d74cc96aca3f9d03f35358b.jpg}
